\section{不足与展望}

\subsection{当前不足}

本文已经完成从数据到演示的完整实验闭环，但仍存在以下限制：

\begin{enumerate}
  \item RUL 泛化仍是主要难点。跨轴承和跨工况差异会导致预测曲线偏移，部分早期模型甚至不能稳定超过 naive baseline，因此 RUL 结论必须保守。
  \item HealthState 和 EarlyFault 是基于退化过程构造的伪标签，不等同于人工拆检或专家标注的真实故障类型。
  \item 当前特征主要是人工统计特征，尚未充分加入 BPFO、BPFI、BSF、FTF 等物理故障频率特征。
  \item tsfresh 在全量长信号上存在资源压力，需要后续设计批处理、降采样或窗口级压缩方案。
  \item 当前训练没有做完整超参数搜索，GRU 结构、学习率、正则化和序列长度仍有优化空间。
  \item 论文复现实验和本文统一三任务框架不是同一实验设置，不能简单横向比较全部指标。
\end{enumerate}

\subsection{未来工作}

后续可以从五个方向改进：

\begin{itemize}
  \item 引入物理故障频率特征，将轴承几何参数和转速信息纳入频域分析。
  \item 比较更强的时序模型，例如 LSTM、TCN、Transformer 或注意力机制模型。
  \item 设计多任务联合学习，让 RUL、HealthState 和 EarlyFault 共享退化表示。
  \item 增加跨工况训练和跨数据集迁移实验，检验模型在更强分布偏移下的鲁棒性。
  \item 引入不确定性估计，避免只给出单点寿命预测而缺少置信度表达。
\end{itemize}

这些方向能够进一步提升方法的工程可信度，也能让本文框架从课程结题实验走向更完整的 PHM 研究系统。
